% =============================================================================
%  BİRİNCİ SAYFA / PAGE ONE  --  doldurulacak form
% =============================================================================
%  Bu dosyadaki her satır ilk sayfada bir yere karşılık gelir. Sadece süslü
%  parantezlerin { } içini değiştirin; komut adlarına dokunmayın.
%
%  Every line in this file corresponds to something on page one. Change only
%  what is inside the braces { }; leave the command names alone.
%
%  Türkçe karakterleri doğrudan yazın: ğ İ ı ş ç ö ü â
%  Yüzde işaretini kaçırın:  \%   ·  ve işaretini:  \&   ·  alt çizgiyi:  \_
% =============================================================================


% ┌───────────────────────────────────────────────────────────────────────────┐
% │  1. DERGİ VE SAYI  /  JOURNAL AND ISSUE                                   │
% │                                                                           │
% │  Bunları bir kez yazın. Üst bant yazısı, sayfa üstü künyesi ve atıf       │
% │  cümlesi bunlardan otomatik üretilir.                                     │
% │  Give these once. The masthead line, the running head and the suggested   │
% │  citation are all built from them.                                        │
% └───────────────────────────────────────────────────────────────────────────┘

\SosaJournal      {Sosyal Bilimler ve Sağlık Bülteni}
\SosaJournalShort {SoSa}
\SosaYear         {2026}
\SosaIssue        {Bahar}{18}        % {sayı adı}{sayı numarası} / {label}{number}
\SosaPages        {60}{70}           % {ilk sayfa}{son sayfa} / {first}{last}

\SosaArticleType  {Özgün Makale}     % Derleme, Olgu Sunumu, Editöre Mektup ...


% ┌───────────────────────────────────────────────────────────────────────────┐
% │  2. BAŞLIK  /  TITLE                                                      │
% │                                                                           │
% │  Türkçe başlık kapak bandının üstüne beyaz olarak yerleşir; iki satırı    │
% │  geçmemelidir. İngilizce başlık bandın hemen altındadır.                  │
% └───────────────────────────────────────────────────────────────────────────┘

\SosaTitleTR{Balıkesir'de Yaşayan Yetişkinlerin Sağlık Hizmetlerinde
             Yapay Zekâ Kullanımına Dair Görüşleri}

\SosaTitleEN{Perspectives of Adults Living in Balıkesir on the Use of
             Artificial Intelligence in Healthcare}

% Sayfa üstü künyesinde kullanılacak kısa başlık. Boş bırakırsanız tam
% Türkçe başlık kullanılır.
% Short title for the running head. Leave it out to use the full title.
% \SosaShortTitle{Yetişkinlerin Sağlık Hizmetlerinde Yapay Zekâ Görüşleri}


% ┌───────────────────────────────────────────────────────────────────────────┐
% │  3. YAZARLAR  /  AUTHORS                                                  │
% │                                                                           │
% │      \SosaAuthor{Ad Soyad}{kurum no}{ORCID}                               │
% │                  ^isim     ^1 veya 1,2   ^boş bırakılabilir               │
% │                                                                           │
% │  Her yazar için bir satır, sırayla. Aralardaki virgüller, üst simgeler   │
% │  ve ORCID rozetleri otomatik eklenir -- siz eklemeyin.                    │
% │  One line per author, in order. Commas, superscripts and ORCID marks are  │
% │  inserted for you -- do not type them.                                    │
% └───────────────────────────────────────────────────────────────────────────┘

\SosaAuthor{Murat Aysin}                {1}{0000-0001-0000-0001}
\SosaAuthor{Sultan Eser}                {2}{0000-0001-0000-0002}
\SosaAuthor{Duygu Lüleci}               {1}{0000-0001-0000-0003}
\SosaAuthor{Emine Ayhan Akman}          {1}{0000-0001-0000-0004}
\SosaAuthor{Çağdaş Sonat}               {3}{0000-0001-0000-0005}
\SosaAuthor{Dursun Yasemin Yayla Keskin}{3}{0000-0001-0000-0006}
\SosaAuthor{Alikemal Başol}             {3}{0000-0001-0000-0007}
\SosaAuthor{Ali Ceylan}                 {4}{0000-0001-0000-0008}


% ┌───────────────────────────────────────────────────────────────────────────┐
% │  4. KURUMLAR  /  AFFILIATIONS                                             │
% │                                                                           │
% │  Yukarıdaki üst simgelerle aynı sırada, her kurum için bir satır.        │
% │  Numaralar otomatik verilir.                                              │
% │  One line per affiliation, in the same order as the superscripts above.   │
% │  They are numbered for you.                                               │
% └───────────────────────────────────────────────────────────────────────────┘

\SosaAffiliation{Dr. Öğr. Üyesi, Balıkesir Üniversitesi, Tıp Fakültesi,
                 Halk Sağlığı Anabilim Dalı. Balıkesir, Türkiye.}
\SosaAffiliation{Prof. Dr., Sağlık Bilimleri Üniversitesi, İzmir Tıp Fakültesi,
                 Halk Sağlığı Anabilim Dalı, İzmir, Türkiye.}
\SosaAffiliation{Arş. Gör. Dr., Balıkesir Üniversitesi, Tıp Fakültesi,
                 Halk Sağlığı Anabilim Dalı. Balıkesir, Türkiye.}
\SosaAffiliation{Prof. Dr., Balıkesir Üniversitesi, Tıp Fakültesi,
                 Halk Sağlığı Anabilim Dalı. Balıkesir, Türkiye.}


% ┌───────────────────────────────────────────────────────────────────────────┐
% │  5. ÖZ  /  TURKISH ABSTRACT                                               │
% │                                                                           │
% │  Alt başlıkları kalın yazın: \textbf{Giriş:} gibi.                        │
% │  Paragrafları BOŞ SATIR ile ayırın.                                       │
% └───────────────────────────────────────────────────────────────────────────┘

\SosaOz{%
  \textbf{Giriş:} Yapay zekâ, sağlık hizmetlerinde tanı, tedavi planlaması,
  hastalık öngörüsü, hasta bakım yönetimi ve klinik karar destek süreçlerinde
  giderek daha fazla kullanım alanı bulmaktadır. Bu teknolojilerin sağlık
  sistemine başarılı biçimde entegre edilebilmesi, toplumun yapay zekâyı
  güvenilir ve faydalı bulmasıyla yakından ilişkilidir. Bu çalışmanın amacı,
  Balıkesir'de yaşayan yetişkinlerin sağlık hizmetlerinde yapay zekâ
  kullanımına dair görüşlerini belirlemektir.

  \textbf{Gereç-Yöntem:} Kesitsel tipteki çalışma, Temmuz 2024'te Balıkesir
  Altıeylül ilçesi merkez mahallelerinde yaşayan 18 yaş ve üzeri yetişkinlerde
  yürütüldü. Çok aşamalı küme örnekleme yöntemiyle 715 kişiye ulaşıldı.
  Veriler yüz yüze anket yöntemiyle Google Forms üzerinden toplandı. Veriler
  SPSS 26.0 programında analiz edildi; kategorik değişkenler için ki-kare;
  sürekli değişkenlerin analizi için t testi kullanıldı. Anlamlılık düzeyi
  p${<}$0,05 kabul edildi.

  \textbf{Bulgular:} Katılımcıların yaş ortalaması 40,3$\pm$15,3 olup
  \%52,9'u kadındı. Katılımcıların \%64,5'i yapay zekânın sağlık
  hizmetlerinde kullanımının faydalı olacağını düşünmekteydi. Katılımcıların
  \%44,8'i yapay zekânın doğru tanılar sunacağını düşünürken, \%57,3'ü
  kişisel verilerin güvenliği ile ilgili endişe duymaktaydı. Tek değişkenli
  analizlerde cinsiyet ve sağlık hizmetlerinde yapay zekânın kullanımında
  kişisel verilerin güvenliği ile ilgili endişe duyma dışında incelenen tüm
  değişkenler ile yapay zekâyı faydalı bulma durumu arasında anlamlı fark
  saptandı.

  \textbf{Sonuç-Öneriler:} Sağlık hizmetlerinde yapay zekâya yönelik fayda
  algısı sosyodemografik özellikler, sağlık durumu ve teletıp deneyimine göre
  değişmektedir. Toplumun farklı gruplarına yönelik güven temelli, anlaşılır
  ve kanıta dayalı bilgilendirme çalışmaları planlanmalıdır.%
}

\SosaKeywordsTR{sağlık hizmeti, yapay zeka, teletıp}


% ┌───────────────────────────────────────────────────────────────────────────┐
% │  6. ABSTRACT  /  ENGLISH ABSTRACT                                         │
% └───────────────────────────────────────────────────────────────────────────┘

\SosaAbstract{%
  \textbf{Introduction:} Artificial intelligence is increasingly being used in
  healthcare, particularly in diagnosis, treatment planning, disease
  prediction, patient care management, and clinical decision support. The
  successful integration of these technologies into the healthcare system is
  closely associated with whether the public perceives artificial intelligence
  as reliable and beneficial. This study aimed to determine the views of
  adults living in Balıkesir regarding the use of artificial intelligence in
  healthcare.

  \textbf{Materials and Methods:} This cross-sectional study was conducted in
  July 2024 among adults aged 18 years and older living in the central
  neighborhoods of Altıeylül district, Balıkesir. A total of 715 participants
  were reached using multistage cluster sampling. Data were collected through
  face-to-face interviews using Google Forms. The data were analyzed using
  SPSS version 26.0. The chi-square test was used for categorical variables,
  and the t-test was used for continuous variables. Statistical significance
  was set at p${<}$0.05.

  \textbf{Results:} The mean age of the participants was 40.3$\pm$15.3 years,
  and 52.9\% were female. Overall, 64.5\% thought that artificial intelligence
  would be beneficial in healthcare. While 44.8\% believed that artificial
  intelligence would provide accurate diagnoses, 57.3\% expressed concerns
  about personal data security. In univariate analyses, all examined
  variables, except sex and concerns about personal data security, differed
  significantly according to whether artificial intelligence was perceived as
  beneficial.

  \textbf{Conclusion and Recommendations:} Perceptions of the benefits of
  artificial intelligence in healthcare vary according to sociodemographic
  characteristics, health status, and telemedicine experience. Clear,
  evidence-based, and trust-oriented informational activities should be
  planned for different groups in society.%
}

\SosaKeywordsEN{healthcare, artificial intelligence, telemedicine.}


% ┌───────────────────────────────────────────────────────────────────────────┐
% │  7. SAYFA ALTI KÜNYESİ  /  COLOPHON AT THE FOOT OF PAGE ONE               │
% └───────────────────────────────────────────────────────────────────────────┘

\SosaReceived            {16.05.2026}
\SosaAccepted            {12.06.2026}
\SosaPublished           {19.06.2026}
\SosaCorrespondingAuthor {Murat Aysin}
\SosaCorrespondingEmail  {murat.aysin@balikesir.edu.tr}
\SosaHandlingEditor      {Nuray Özgülnar}

% Atıf cümlesindeki yazar listesi. Gerisi (yıl, başlık, dergi, sayfalar)
% yukarıdaki bilgilerden otomatik tamamlanır.
% Just the author list for the citation. The year, title, journal and pages
% are completed from the fields above.
\SosaCitationAuthors{Aysin, M., Eser, S., Lüleci, D., Ayhan Akman, E.,
  Sonat, Ç., Yayla Keskin, D. Y., Başol A., \& Ceylan A.}

% Atıf cümlesini tümüyle kendiniz yazmak isterseniz bunu açın; yukarıdakini
% geçersiz kılar. To write the whole citation by hand, uncomment this:
% \SosaCitation{Soyad, A. (2026). Başlık. \textit{Dergi, SAYI(18),} 60-70}


% ┌───────────────────────────────────────────────────────────────────────────┐
% │  8. İSTEĞE BAĞLI NOT  /  OPTIONAL NOTE                                    │
% │                                                                           │
% │  Özetlerin altında görünür: kongre bildirisi, destek, tez bilgisi vb.     │
% │  Gerekmiyorsa satırı yorumda bırakın.                                     │
% │  Appears under the abstracts. Leave it commented out if not needed.       │
% └───────────────────────────────────────────────────────────────────────────┘

% \SosaNote{Bu makalede yer alan verilerin bir kısmı 27.12.2025 tarihinde
%   Konya'da düzenlenen Uluslararası Selçuk Bilimsel Araştırmalar
%   Kongresi'nde sunulmuştur.}
